\documentclass[11pt]{article}
\usepackage[a4paper,margin=1in]{geometry}
\usepackage{fontspec}
\usepackage[boldfont=false]{xeCJK}
\setCJKmainfont{FandolSong-Regular.otf}
\usepackage{amsmath}
\usepackage{amssymb}
\usepackage{array}
\usepackage{booktabs}
\usepackage{graphicx}
\usepackage{adjustbox}
\usepackage{enumitem}
\setlist[itemize]{topsep=2pt,partopsep=0pt,parsep=0pt,itemsep=2pt,leftmargin=1.4em}
\setlist[enumerate]{topsep=2pt,partopsep=0pt,parsep=0pt,itemsep=2pt,leftmargin=1.6em}
\usepackage{parskip}
\usepackage{fvextra}
\fvset{breaklines=true,breakanywhere=true,fontsize=\small}
\setlength{\parindent}{0pt}

\begin{document}

\begin{center}
  {\LARGE\bfseries International economic issues (A Level)}\\[0.35em]
  {\large A Level Economics}
\end{center}
\vspace{0.6em}

This final topic covers trade imbalances, exchange rates in depth, and how countries develop and connect.

\section*{Correcting a balance of payments imbalance}

A country with a large \textbf{balance of payments} 国际收支 problem — usually a deficit on the \textbf{current account} 经常账户 — can use two kinds of policy.

\begin{itemize}
\item \textbf{expenditure-reducing} 支出减少 policies cut total demand (higher taxes or interest rates). With less spending, people buy fewer imports. But growth slows and unemployment may rise.

\item \textbf{expenditure-switching} 支出转换 policies move spending away from imports towards home goods. The main tool is a lower \textbf{exchange rate} 汇率 — a \textbf{depreciation} 贬值 under floating rates, or a \textbf{devaluation} 法定贬值 under a fixed system — which makes exports cheaper and imports dearer.

\end{itemize}

\subsection*{The Marshall–Lerner condition and the J-curve}

A lower exchange rate only improves the current account if buyers respond enough. The \textbf{Marshall-Lerner condition} 马歇尔勒纳条件 says the sum of the price elasticities of demand for exports and imports must be greater than one:

\[
PED_{\text{exports}} + PED_{\text{imports}} > 1
\]

Even when this holds, the gain takes time. The \textbf{J-curve effect} J曲线效应 describes what happens: just after a depreciation, the current account first gets \textbf{worse} (imports cost more straight away, but trade volumes are slow to change), then later improves as buyers adjust — tracing the shape of a letter J.

\section*{Exchange rates}

An \textbf{exchange rate} is the price of one currency in terms of another. There is more than one way to measure it:

\begin{itemize}
\item the \textbf{nominal exchange rate} 名义汇率 is the everyday market rate between two currencies.

\item the \textbf{real exchange rate} 实际汇率 adjusts the nominal rate for differences in prices between the countries, so it shows true competitiveness.

\item the \textbf{trade-weighted exchange rate} 贸易加权汇率, also called the \textbf{effective exchange rate} 有效汇率, is an average of a currency's value against all its trading partners, weighted by how much trade is done with each.

\end{itemize}

\subsection*{The three systems}

\begin{itemize}
\item a \textbf{floating exchange rate} 浮动汇率 is set by the demand for and supply of the currency.

\item a \textbf{fixed exchange rate} 固定汇率 is held at a set value by the central bank.

\item a \textbf{managed exchange rate} 管理浮动汇率 mostly floats, with the central bank stepping in at times.

\end{itemize}

Under floating rates, a rise in value is an \textbf{appreciation} 升值 and a fall is a \textbf{depreciation}. Under a fixed system, a deliberate rise is a \textbf{revaluation} 法定升值 and a deliberate fall is a \textbf{devaluation}. A weaker currency tends to raise exports and inflation; a stronger currency does the opposite.

\section*{Economic development}

It is important not to confuse two ideas:

\begin{itemize}
\item \textbf{economic growth} 经济增长 is simply a rise in real output.

\item \textbf{economic development} 经济发展 is wider. It is about people's \textbf{standard of living} 生活水平 — their health, education, and freedom of choice, not just their income.

\end{itemize}

A country can grow without developing much (if the extra income goes to a few people). Development usually means rising incomes \textbf{plus} better health, more schooling, and less poverty.

\subsection*{Measuring development}

The best-known measure is the \textbf{Human Development Index} 人类发展指数 (HDI). It combines three things into one number between 0 and 1:

\begin{itemize}
\item income (GNI per person),

\item health (\textbf{life expectancy} 预期寿命 at birth),

\item education (years of schooling and adult \textbf{literacy} 识字率).

\end{itemize}

Other useful signs of development include access to clean water, the number of doctors per person, and mobile-phone and internet use.

\section*{Comparing countries at different levels of development}

We often compare a \textbf{developed country} 发达国家 (high income, like Germany) with a \textbf{developing country} 发展中国家 (lower income, like Kenya). They differ in clear ways:

\begin{center}
\adjustbox{max width=\linewidth}{%
\begin{tabular}{lll}
\toprule
\textbf{Feature} & \textbf{Developed} & \textbf{Developing} \\
\midrule
income per person & high & low \\
\textbf{birth rate} 出生率 and population growth & low & often high \\
life expectancy & high & lower \\
share of workers in farming & small & large \\
\bottomrule
\end{tabular}}
\end{center}

The \textbf{employment structure} 就业结构 also differs. As a country develops, workers move out of the \textbf{primary sector} 第一产业 (farming and mining), into the \textbf{secondary sector} 第二产业 (making goods), and later into the \textbf{tertiary sector} 第三产业 (services).

Many things shape development: savings and investment, the quality of education and health, infrastructure, stable government and law, and access to world markets.

\section*{Relationships between countries}

Richer and poorer countries are linked in several ways, each with good and bad sides.

\begin{itemize}
\item \textbf{trade} can raise incomes, but many developing countries rely on a few primary exports whose prices swing sharply.

\item foreign \textbf{aid} 援助 can fund schools and roads, but it may create dependency or be wasted.

\item \textbf{debt} 债务 owed to other countries can pay for development, but repaying it can drain a poor country's budget.

\item \textbf{foreign direct investment} 外国直接投资 (FDI) by a \textbf{multinational corporation} 跨国公司 brings money, jobs and technology, but profits may flow back abroad and workers may be poorly treated.

\item \textbf{migration} 移民 of workers can ease unemployment and bring in \textbf{remittances} 侨汇 (money sent home), but it can also drain a country of its most skilled people.

\end{itemize}

\section*{Globalisation}

\textbf{globalisation} 全球化 is the growing links between countries — in trade, money, technology and people — so that the world acts more like one market. Its causes include cheaper transport and communication, \textbf{free trade} 自由贸易 agreements, and the spread of multinational firms.

Countries often join a \textbf{trade bloc} 贸易集团 — a group that trades freely among themselves (like the European Union).

\subsection*{Benefits and costs}

\begin{center}
\adjustbox{max width=\linewidth}{%
\begin{tabular}{ll}
\toprule
\textbf{Benefits} & \textbf{Costs} \\
\midrule
lower prices and more choice & some industries and jobs are lost to cheaper countries \\
faster growth and technology transfer & the gap between rich and poor can widen \\
larger markets for firms & more pollution and use of resources \\
poorer countries can industrialise & powerful multinationals may avoid tax and dodge \textbf{protectionism} 贸易保护主义 rules \\
\bottomrule
\end{tabular}}
\end{center}

Globalisation helps developed and developing economies in different ways, but the gains are not shared evenly, which is why it remains so argued about.


\end{document}
